\documentclass[12pt,a4paper]{article}
\usepackage[margin=25mm, headheight=15pt, headsep=10pt]{geometry}
\usepackage{fontspec}
\usepackage[table]{xcolor} % Note: table option is required for rowcolors
\usepackage{titlesec}
\usepackage{tocloft}
\usepackage{fancyhdr}
\usepackage{booktabs}
\usepackage{tcolorbox}
\tcbuselibrary{skins, breakable}
\usepackage[colorlinks=true, linkcolor=emerald, urlcolor=emerald, bookmarks=true]{hyperref}

\definecolor{emerald}{RGB}{0,102,51}
\definecolor{darkred}{RGB}{139,0,0}
\definecolor{lightgray}{RGB}{245,245,245}

\usepackage[nil]{babel}
\babelprovide[import, main]{bengali}
\babelprovide[import]{arabic}
\babelfont{rm}[Renderer=HarfBuzz,Scale=1.0]{Noto Serif Bengali}
\babelfont{sf}[Renderer=HarfBuzz]{Noto Sans Bengali}
\babelfont{tt}[Renderer=HarfBuzz]{Noto Sans Bengali}
\babelfont[arabic]{rm}[Renderer=HarfBuzz,Scale=1.1]{Noto Naskh Arabic}

% Section styling
\titleformat{\section}{\Large\bfseries\color{emerald}}{\thesection.}{0.5em}{}
\titleformat{\subsection}{\large\bfseries\color{emerald!85!black}}{\thesubsection.}{0.5em}{}
\titleformat{\subsubsection}{\normalsize\bfseries\color{emerald!70!black}}{\thesubsubsection.}{0.5em}{}
\titlespacing*{\section}{0pt}{18pt}{8pt}

% Fancy Header/Footer setup
\pagestyle{fancy}
\fancyhf{} % clear all header and footer fields
\fancyhead[L]{\color{emerald}\small\textbf{\nouppercase{\leftmark}}}
\fancyfoot[L]{\scriptsize\color{black!50}{\lat{AI}}-সংকলিত, আলেম-যাচাইকৃত নয়}
\fancyfoot[C]{\color{emerald!70!black}\thepage}
\renewcommand{\headrulewidth}{0.4pt}
\renewcommand{\headrule}{\hbox to\headwidth{\color{emerald}\leaders\hrule height \headrulewidth\hfill}}
\renewcommand{\footrulewidth}{0pt}

% Custom boxes
% 1. Ayat/Hadith Box
\newtcolorbox{ayatbox}[1][]{
    enhanced, breakable, 
    colback=emerald!5, 
    colframe=emerald, 
    boxrule=0.5pt, 
    arc=3pt, 
    left=10pt, right=10pt, top=10pt, bottom=10pt,
    #1
}

% 2. Quote/Translation Box
\newtcolorbox{quotebox}[1][]{
    enhanced, breakable,
    frame hidden,
    colback=lightgray,
    borderline west={3pt}{0pt}{emerald},
    left=15pt, right=10pt, top=10pt, bottom=10pt,
    sharp corners,
    #1
}

% 3. AI-generated notice box
\newtcolorbox{warnbox}[1][]{
    enhanced, breakable,
    colback=red!4,
    colframe=darkred,
    boxrule=0.8pt,
    arc=3pt,
    left=10pt, right=10pt, top=10pt, bottom=10pt,
    #1
}

% Commands
\newcommand{\arab}[1]{\foreignlanguage{arabic}{#1}}
\newcommand{\kib}[1]{\textcolor{emerald}{\textbf{#1}}}

% Latin (English) fallback: Noto Serif Bengali has NO Latin glyphs
\newfontfamily\latfont[Renderer=HarfBuzz]{Noto Serif}
\newcommand{\lat}[1]{{\latfont #1}}

\setlength{\parskip}{5pt}
\setlength{\parindent}{0pt}

% Fix for missing bullet in Noto Serif Bengali
\renewcommand{\labelitemi}{$\bullet$}



\begin{document}
\begin{center}{\Large\bfseries\color{emerald} পার্ট ৬ — বিতর্কিত মাসআলা: মাযহাব-তুলনা, দলিল ও সমাধান}\end{center}
\vspace{1em}
\section*{পার্ট ৬ — বিতর্কিত মাসআলা: মাযহাব-তুলনা, দলিল ও সমাধান}

\begin{warnbox}
 \textbf{গুরুত্বপূর্ণ নোটিশ (\lat{Important} \lat{Notice}):} এই কন্টেন্টটি \lat{AI} (কৃত্রিম বুদ্ধিমত্তা) দিয়ে প্রস্তুত ও সংকলিত। \lat{AI} কঠোর সূত্র-মান নীতি মেনে শুধু নির্ভরযোগ্য উৎস থেকে তথ্য সংগ্রহ করেছে — কেবল সহীহ/হাসান হাদিস ও সহীহ সনদের আসার রাখা হয়েছে, দুর্বল সনদ বাদ দেওয়া হয়েছে। তবে এটি কোনো আলেম বা মুহাদ্দিস কর্তৃক যাচাইকৃত নয়।
\end{warnbox}

\begin{quote}
\textbf{সম্পর্কিত ফাইল:} পার্ট ২-৫ঘ (মাসআলা), পার্ট ৭ক-৭খ (ভুল ও ভঙ্গকারী)।
\end{quote}

\begin{quote}
\textbf{যে প্রশ্নের উত্তর এই অংশে:} নামাজের যে বিষয়গুলোতে মুসলিমরা বিভক্ত (হাত তোলা, আমিন, ফাতিহা, সূরা-মিলানো, তারাবীহ-সংখ্যা...) — সেগুলোর \textbf{দলিল ও মাযহাব-অবস্থান} কী? কেন মতভেদ হলো? একজন সাধারণ মুসল্লি কীভাবে নিরপেক্ষভাবে বোঝেন?
\end{quote}

\medskip\hrule\medskip

\subsection*{১. ভূমিকা: কেন মতভেদ (\lat{Ikhtilaf}) স্বাভাবিক}

নবীজির (সা.) যুগে নামাজের মূল কাঠামো এক — তবে \textbf{সূক্ষ্ম বিস্তারিত} (কোন হাত তোলা হবে, আমিন জোরে/আস্তে, ফাতিহা পড়বে কিনা) নিয়ে সাহাবা ও তাবেয়িনদের মধ্যে ভিন্ন ভিন্ন বর্ণনা এবং ভিন্ন ভিন্ন আমল ছিল। ফকিহরা (মুহাদ্দিস-ফকিহ) দলিল ও নীতি দিয়ে \textbf{বৈধ ইখতিলাফের} মধ্যে নিজস্ব মত গড়েছেন।

\textbf{কেন পার্থক্য হয়? (৪ মূল কারণ):}
\begin{enumerate}
\item \textbf{হাদিস পৌঁছানো না (বুলুঘ):} কিছু হাদিস ইমামের কাছে পৌঁছেনি।
\item \textbf{হাদিসের গ্রহণ-মান (নাসিখ-মানসুখ):} কেউ বলেন রহিত হয়েছে, কেউ বলেন বহাল।
\item \textbf{কুরআনের ব্যাখ্যা (তাফসির-বৈচিত্র্য):} এক আয়াত একাধিক অর্থে।
\item \textbf{কিয়াস-ও-ইসতিসহাব:} দলিল না পেলে কিয়াস (\lat{analogy}) বা মূল-নীতি প্রযোজ্য।
\end{enumerate}
\begin{quote}
\textbf{মূলনীতি:} \textbf{বৈধ ইখতিলাফে কোনো পক্ষকে "ভুল" বলা যায় না} — প্রতিটি মাযহাবের নিজস্ব দলিল আছে; কেউ কুরআন-সুন্নাহ বাদ দেয়নি। সাধারণ মুসল্লির করণীয়: \textbf{তার মাযহাব/ইমামের অনুসরণ} — এবং অন্যের মাযহাবের বৈধতা স্বীকার।
\end{quote}

\medskip\hrule\medskip

\subsection*{২. সবচেয়ে আলোচিত ১০টি বিতর্কিত মাসআলা}

\subsubsection*{২.১ রাফউল ইয়াদাইন (হাত তোলা)}

\begin{itemize}
\item \textbf{দলিল (সহীহ বুখারি ৭৩৫):} ইবনে উমার (রাঃ) — \lat{“}নবীজি (সা.) তাকবিরে তাহরিমায়, রুকুতে যাওয়ার সময় ও রুকু থেকে উঠার সময় দুই হাত কাঁধ বরাবর তুলতেন।\lat{”}
\item \textbf{হানাফি:} শুধু তাকবিরে তাহরিমায় হাত তোলা সুন্নত; বাকি জায়গায় না — কারণ তাদের মতে \textbf{রহিত (মানসুখ)} হয়েছে (ইবনে মাসউদ (রাঃ)-এর বর্ণনায় হাত তোলার বাদ-অংশ)।
\item \textbf{শাফেয়ি-হাম্বলি:} প্রতিটি তাকবিরে হাত তোলা সুন্নত।
\item \textbf{মালেকি:} তাকবিরে তাহরিমায় উত্তম।
\item \textbf{সারমর্ম:} হাত তোলা \textbf{সুন্নত}, ফরজ-ওয়াজিব নয় — কোনো মাযহাবই এটা ফরজ বলে না। পার্থক্য শুধু \lat{“}কোথায়\lat{”}।
\end{itemize}
\subsubsection*{২.২ আমিন — জোরে না আস্তে?}

\begin{itemize}
\item \textbf{দলিল:} \lat{“}যখন ইমাম আমিন বলেন, তোমরাও আমিন বলো\lat{”} (সহীহ বুখারি ৭৮০)।
\item \textbf{জোরে (জাহরি) আমিন:} হানাফি-হাম্বলি (প্রসিদ্ধ) — জোরে; হাম্বলি মতেও জোরে।
\item \textbf{আস্তে (সিররি):} মালেকি-শাফেয়ি — আস্তে (শাফেয়ি: জোরে নয়)।
\item \textbf{সারমর্ম:} আমিন \textbf{সব মাযহাবেই আছে} — পার্থক্য শুধু জোরে-আস্তে; দলিল দুইভাবে বর্ণিত, নবীজি (সা.) দুইভাবে করেছেন (রিপোর্ট দুই রকম)।
\end{itemize}
\subsubsection*{২.৩ ইমামের পেছনে ফাতিহা পড়া}

\begin{itemize}
\item \textbf{দলিল (সহীহ বুখারি ৭৫২):} \lat{“}যে নামাজ পড়ল অথচ ফাতিহা পড়ল না, তার নামাজ অপূর্ণ।\lat{”}
\item \textbf{দলিল (সহীহ মুসলিম ৪০৪):} \lat{“}যে ইমাম রুকু করল (পেছনে ফাতিহা পড়ে) — এতে কিছু যায়-আসে না\lat{”} (আবু হুরাইরা থেকে — পরোক্ষ)।
\item \textbf{শাফেয়ি-হাম্বলি:} ইমামের পেছনে ফাতিহা \textbf{ফরজ} (জোরে-নিরবে দুই-ই)।
\item \textbf{মালেকি:} ফাতিহা পড়বে (ফরজ-এক মত); ইমামের নিরবে নামাজে পড়া।
\item \textbf{হানাফি:} ইমামের পেছনে \textbf{ফাতিহা নয়} — ইমামের কিরাআতই যথেষ্ট (দলিল: সূরা আল-আ'রাফ ৭:২০৪ \lat{“}যখন কুরআন পাঠ করা হয়, চুপ থাকো/শোনো\lat{”} + বুখারি ৭২২-এর বর্ণনা)।
\item \textbf{সারমর্ম:} এটিই সবচেয়ে গভীর ইখতিলাফ — দুই পক্ষেরই শক্ত দলিল; \textbf{মসজিদে ইমামের অনুসরণই} বাস্তব সমাধান; নিজ মাযহাব অনুযায়ী আমল।
\end{itemize}
\subsubsection*{২.৪ ফাতিহার পরে সূরা মিলানো — প্রতিটি রাকাতে?}

\begin{itemize}
\item \textbf{হানাফি:} প্রথম দুই রাকাতে \textbf{ফাতিহা+সূরা} (ওয়াজিব); ৩য়-৪র্থ রাকাতে শুধু ফাতিহা (সূরা মিলানো সুন্নত)।
\item \textbf{জমহুর:} প্রতিটি রাকাতে ফাতিহা; সূরা মিলানো সুন্নত।
\item \textbf{দলিল:} নবীজি (সা.) প্রথম দুই রাকাতে জোরে, শেষ দুই-এ আস্তে কিরাআত করেছেন (সহীহ বুখারি ৭৫৯) — সূরা মিলানোর স্পষ্ট বর্ণনা সব রাকাতে নেই — তাই মাযহাবভেদে ভিন্ন।
\end{itemize}
\subsubsection*{২.৫ বিসমিল্লাহ জোরে বলা}

\begin{itemize}
\item \textbf{হানাফি:} আস্তে (নিরবে) — ফাতিহা শুরুর আগে।
\item \textbf{শাফেয়ি:} \textbf{জোরে} (জাহরি নামাজে) — দলিল: সহীহ-হাসান বর্ণনায় নবীজি জোরে বিসমিল্লাহ পড়েছেন (সহীহ বুখারি ৭৪৩-এর বর্ণনা কিছুটা ভিন্ন)।
\item \textbf{মালেকি-হাম্বলি:} আস্তে (মালেকি: শুধু ফাতিহার আগে)।
\item \textbf{সারমর্ম:} বিসমিল্লাহ তিলাওয়াত \textbf{সবাই করে} — পার্থক্য শুধু জোরে-আস্তে।
\end{itemize}
\subsubsection*{২.৬ তাশাহহুদের আঙুল (ইশারা)}

\begin{itemize}
\item \textbf{সহীহ হাদিস (সহীহ মুসলিম ৫৮০):} নবীজি (সা.) তাশাহহুদে তর্জনী দিয়ে ইশারা করেছেন — \textbf{আঙুল ওঠানো প্রমাণিত সুন্নত}।
\item \textbf{হানাফি:} \lat{“}লা-ইলাহা\lat{”}-য় আঙুল ওঠানো; চোখ ইশারায়।
\item \textbf{শাফেয়ি-হাম্বলি:} আঙুল ওঠানো।
\item \textbf{মালেকি:} \textbf{আঙুল না ওঠানো} — সিজদার মতো বিছিয়ে রাখা (মালেকি দলিল: নবীজির আমলের ভিন্ন বর্ণনা)।
\item \textbf{সারমর্ম:} আঙুল ইশারা \textbf{সুন্নত-সংক্রান্ত} — কেউ ফরজ বলে না; মালেকি মত আলাদা আমল হলেও বৈধ।
\end{itemize}
\subsubsection*{২.৭ কুনুত — নাযিলা (বিপদে) ও বিতরে}

\begin{itemize}
\item \textbf{কুনুত নাযিলা:} বিপদ-সংকটে নামাজে দুআ — নবীজি (সা.) করেছেন (সহীহ বুখারি ১০০২-১০০৩; সহীহ মুসলিম ৬৭৬)। \textbf{সব মাযহাবে} জায়েয — পার্থক্য কখন: হানাফি-মালেকি শুধু বিপদে; শাফেয়ি-হাম্বলি রমজানের দ্বিতীয় অর্ধে বিতরে নিয়মিত।
\item \textbf{বিতরের কুনুত:} হানাফি — ৩য় রাকাতে ওয়াজিব; জমহুর — সুন্নত (দুআ-কুনুত হাদিস — সুনানে আবু দাউদ ১৪২৫)।
\end{itemize}
\subsubsection*{২.৮ তারাবীহ-রাকাত: ৮ না ২০?}

\begin{itemize}
\item \textbf{দলিল (২০):} উমার (রাঃ)-এর যুগে ২০ রাকাত জামাতের সহীহ-হাসান সনদ (আবদুর রহমান ইবনে আবদিল ক্বারী — মালিকের মুওয়াত্তা/সহীহ)।
\item \textbf{দলিল (৮):} আয়েশা (রাঃ) — নবীজি (সা.) রমজানে ১১ রাকাত (৮ তাহাজ্জুদ+৩ বিতর) — (সহীহ বুখারি ২০১২-এর বর্ণনা)।
\item \textbf{মাযহাব:} হানাফি-শাফেয়ি-হাম্বলি: \textbf{২০}; মালেকি: \textbf{৩৬} (মদিনা-আমল); কিছু সালাফী-মত: ৮-২০ দুই-ই বৈধ।
\item \textbf{সারমর্ম:} \textbf{উমার (রাঃ) ও নবীজি-যুগের আলাদা আমল} — দুই দলিলই বৈধ; রাকাত-সংখ্যা \textbf{ইখতিলাফি} — মূল বিষয় জামাতে কিয়াম-তিলাওয়াত। কোনো সংখ্যাকে হারাম/ভুল বলা যায় না।
\end{itemize}
\subsubsection*{২.৯ বিতর: ১ না ৩? কুনুত?}

\begin{itemize}
\item \textbf{দলিল (১ রাকাত):} নবীজি (সা.) মাঝে মাঝে ১ রাকাত বিতর পড়েছেন (সহীহ বুখারি ৯৯০)।
\item \textbf{দলিল (৩ রাকাত):} নবীজি (সা.) প্রায়ই ৩ রাকাত বিতর — সালাম না ফিরে (সহীহ বুখারি ৯৯৩-এর বর্ণনা; সহীহ মুসলিম ৭৫২)।
\item \textbf{হানাফি:} ৩ রাকাত — এক সালামে, কুনুতসহ (ওয়াজিব)।
\item \textbf{জমহুর:} ৩ রাকাত \textbf{বা} ১ রাকাত — দুই-ই জায়েয; কুনুত সুন্নত।
\end{itemize}
\subsubsection*{২.১০ মোজার উপর মাসেহ-এর সময়সীমা}

\begin{itemize}
\item হানাফি-শাফেয়ি-হাম্বলি: মুকিম ১ দিন-১ রাত, মুসাফির ৩ দিন-৩ রাত (সহীহ মুসলিম ২৭৬)।
\item \textbf{মালেকি:} \textbf{কোনো সময়সীমা নেই} — মাসেহ ততক্ষণ মোজা পবিত্র থাকা পর্যন্ত; দলিল: মোজা পরা অবস্থায় হাদাস ছাড়া ওযু পুনরায়। (বৈধ ইখতিলাফ)
\end{itemize}
\medskip\hrule\medskip

\subsection*{৩. কীভাবে মীমাংসা করবেন? (\lat{Practical} \lat{Guide})}

\begin{enumerate}
\item \textbf{ইমাম-মুসল্লির বাস্তব নিয়ম:} মসজিদে ইমামের আমলই অনুসরণ করুন — নিজের মাযহাব ভিন্ন হলে ইমামের অনুসরণ \textbf{ভালো} (জামাতের একতা রক্ষা)।
\item \textbf{নিজের মাযহাবের জ্ঞান নিন:} হানাফি হলে হানাফি-ফিকহের গ্রন্থ (যেমন নূরুল ইদাহ, ফিকহুল ইসলামি) পড়ুন; শাফেয়ি হলে সাবানুন-নাজাত ইত্যাদি।
\item \textbf{ইখতিলাফি মাসআলায় অন্যকে "ভুল" বলা নয়:} হাত তোলা-না তোলা, আমিন জোরে-আস্তে — কাউকে গুমরাহ বলা হারাম-সদৃশ। \textbf{ইখতিলাফ রহমত} — কিন্তু দলিল-জ্ঞান ছাড়া নিজের মতই শ্রেষ্ঠ বলে রায় দেবেন না।
\item \textbf{সন্দেহ হলে:} নির্ভরযোগ্য আলেম-মুফতির কাছে প্রশ্ন — নিজে নিজে ফতোয়া-সিদ্ধান্ত নেবেন না।
\end{enumerate}
\medskip\hrule\medskip

\subsection*{৪. তুলনা-সারণি (সংক্ষিপ্ত)}

\begin{table}[h]\centering\small
\begin{tabular}{lllll}
মাসআলা & হানাফি & শাফেয়ি & মালেকি & হাম্বলি \\
\hline
রাফউল ইয়াদাইন & শুধু তাকবিরে & প্রতিটি তাকবিরে & তাকবিরে & প্রতিটি তাকবিরে \\
আমিন & জোরে & আস্তে & আস্তে & জোরে \\
পেছনে ফাতিহা & না (ইমামের কিরাআত) & হ্যাঁ (ফরজ) & হ্যাঁ & হ্যাঁ (ফরজ) \\
বিসমিল্লাহ & আস্তে & জোরে & আস্তে & আস্তে \\
তাশাহহুদে আঙুল & ওঠায় & ওঠায় & না & ওঠায় \\
তারাবীহ & ২০ & ২০ & ৩৬ & ২০ \\
বিতর & ৩+কুনুত & ১/৩+কুনুত & ১/৩ & ১/৩+কুনুত \\
\hline
\end{tabular}
\end{table}

\medskip\hrule\medskip

\subsection*{৫. রেফারেন্স}

\begin{itemize}
\item সহীহ বুখারি — ৭৩৫, ৭৪৩, ৭৫২, ৭৫৯, ৭৮০, ৯৯০, ৯৯৩, ১০০২-১০০৩, ২০১২
\item সহীহ মুসলিম — ২৭৬, ৪০৪, ৫৮০, ৬৭৬, ৭৫২
\item সুনানে আবু দাউদ — ১৪২৫ (কুনুত)
\item সুনানে তিরমিজি — ৪১৪
\item মুওয়াত্তা ইমাম মালিক (উমার-যুগের ২০ রাকাত তারাবীহ)
\item আল-ফিকহুল ইসলামি ওয়া আদিল্লাতুহু (আল-যুহাইলি) — ৪-মাযহাব তুলনার প্রধান রেফারেন্স
\end{itemize}
\begin{quote}
\textbf{দ্রষ্টব্য:} এই অংশে কোনো মাযহাবকে শ্রেষ্ঠ বলার চেষ্টা নেই — \textbf{দলিল-তুলনা ও ইখতিলাফ-বোঝার} জন্য। প্রতিটি মাসআলার শক্ত-দুর্বল দিক পার্ট ৭ক-৭খ-এ বাস্তব প্রয়োগে আরও দেখা যাবে।
\end{quote}

\end{document}
